\section{Introduction}\label{sec:introduction}

In metric \(k\)-median, a finite set \(J\) of clients and a finite set \(F\)
of candidate facilities lie in a common metric space. Given an integer \(k\),
one chooses at most \(k\) facilities to minimize the sum of the distances
from clients to their nearest chosen facility. For a nonempty set
\(S\subseteq F\), write
\[
 d(j,S)=\min_{i\in S}d(j,i),\qquad
 \cost(S)=\sum_{j\in J}d(j,S),\qquad
 \OPT_k=\min_{\substack{S\subseteq F\\1\le|S|\le k}}\cost(S).
\]
The input distances are rational and encoded in binary, and
\(1\le k\le |F|\). Both \(k\) and the metric are unrestricted.
Candidate facilities are specified as part of the input; we do not
require \(F=J\). The empty-client instance has cost zero and can be
handled immediately.

The facility budget is a central difficulty: a solution with small
connection cost can be much easier to obtain when extra facilities are
allowed. Our approach first permits a constant additive surplus, makes
that approximation deterministic, and then removes the surplus using an
established reduction.

\begin{theorem}\label{thm:main}
For every fixed \(\varepsilon>0\), there is a deterministic algorithm which,
given a finite rational metric \(k\)-median instance with \(1\le k\le |F|\),
returns \(S\subseteq F\), \(|S|\le k\), such that
\[
 \sum_{j\in J}d(j,S)\le
 \left(1+\frac2e+\varepsilon\right)\OPT_k.
\]
Its running time is polynomial in the binary encoding length of the instance.
The polynomial may depend on \(\varepsilon\).
\end{theorem}

Hardness below \(1+2/e\) in this facility model follows from the
Max-\(k\)-Coverage gap: covered clients have distance one and uncovered
clients have distance at least three. We use the formulation under \(P\ne NP\)
recorded by Anand and Lee~\cite[Sections~1 and~1.1]{AnandLee2024}, based on
Feige's perfect-completeness coverage gap~\cite[Theorem~5.3]{Feige1998}.
For an explicit metric completion, represent sets by facilities and
universe elements by a disjoint client set. Give distinct points within
either part distance two, and give a facility--client pair distance one
for incidence and three otherwise. This is a metric: a same-part edge
of length two is bounded by \(1+1\), and a cross-part edge of length at
most three is bounded by \(2+1\). For a universe of size \(m\), the
connection cost is exactly \(3m-2c\) when the chosen sets cover \(c\)
elements; a perfect cover has cost \(m\). Thus a gap between
covering every element and covering at most a
\(1-1/e+\delta\) fraction gives the lower bound
\(1+2/e-2\delta\), for every fixed positive \(\delta\).
Together with Theorem~\ref{thm:main},
this proves the following consequence.

\begin{corollary}\label{cor:threshold}
If \(P\ne NP\), the infimum of deterministic polynomial-time approximation
factors for metric \(k\)-median with specified candidate facilities is \(1+2/e\).
\end{corollary}

The infimum in Corollary~\ref{cor:threshold} does not assert attainment of the
endpoint factor. The hardness assumption and candidate-facility model are part
of the statement. In particular, the corresponding lower-bound question with
the restriction \(F=J\) is different~\cite{AnandLee2024}.

\subsection{Historical context and method}

Metric \(k\)-median has been a testing ground for LP rounding,
primal-dual algorithms, and local search. Charikar, Guha, Tardos and
Shmoys gave the first constant-factor approximation through LP
rounding~\cite{CharikarGuhaTardosShmoys2002}. Jain and
Vazirani connected the cardinality constraint with metric facility
location through Lagrangian relaxation and the primal-dual
method~\cite{JainVazirani2001}. The combinatorial improvements of
Charikar and Guha~\cite{CharikarGuha1999} and the greedy dual-fitting
approach of Jain et al.~\cite{JainEtAl2003} developed this connection
further. A different route, the bounded-swap local search of Arya et
al.~\cite{Arya2004}, achieves factor \(3+\varepsilon\).

Li and Svensson~\cite{LiSvensson2016} made constant-additive
pseudo-approximations a route past the factor-three barrier: a constant
number of extra facilities can be removed at an arbitrarily small
additional approximation loss. Subsequent bi-point rounding, including
the corrected \(2.675+\varepsilon\) guarantee of Byrka et
al.~\cite{ByrkaEtAl2017}, refined this approach. Gowda et
al.~\cite{GowdaEtAl2023} obtained \(2.613\) in their SODA 2023 version.
Cohen-Addad et al.~\cite{CGLSS2026} reached
\(2+\varepsilon\), first announced in 2025. Byrka et
al.~\cite{ByrkaEtAl2026} subsequently developed a graph-based iterative
rounding framework that also gives \(2+\varepsilon\) for metric
\(k\)-median and extends to other clustering objectives.
Theorem~\ref{thm:main} reaches every fixed factor above the
\(1+2/e\) hardness boundary in the general candidate-facility model.

Our construction follows the iterative graph framework, with the
cost analysis replaced by a directed comparison inequality and a
potential that retains some surviving assignments after an opening.
The earlier framework's general transfer theorem uses a parameter at
least \(2\) in the metric \(k\)-median case
\cite[Theorem~6.1]{ByrkaEtAl2026}; it does not supply the smaller factor
proved here. We establish the grid preparation, physical updates,
comparison, and deterministic implementation together.

A separate companion combines anchor recovery with bounded-price payment
accounting to give an independent randomized approximation below factor
\(2\) while respecting the facility budget~\cite[Theorem~1.2]{OpenAIRecovery2026}.

Two other methodological antecedents are useful to distinguish.
Splitting fractional facilities and organizing assignments around
separated representatives and bundles are standard LP-rounding
techniques; Charikar and Li~\cite{CharikarLi2012} provide a close
predecessor. The pair updates used in our preparation belong to the
dependent-rounding approach of Srinivasan~\cite{Srinivasan2001} and
Gandhi et al.~\cite{GandhiEtAl2006}, with weighted variants studied by
Byrka et al.~\cite{ByrkaEtAl2017}. We prove the particular moment bounds
and coupled supplier construction needed here. Likewise, reducing the
support of a distribution while matching moments is a standard
Carath\'eodory argument; see Bayer and Teichmann~\cite{BayerTeichmann2005}
for the finite-cubature perspective. Our implementation specifies which
moments suffice throughout the physical process and proves the required
rational encoding bounds.

\subsection{Proof overview}

The proof separates the preparation of fractional openings, their
iterative rounding, and the removal of randomness and surplus. The
intermediate guarantee is measured against the standard assignment LP:
its variable \(y_i\) is the amount of opening at facility \(i\), and each
client receives a unit of this opening mass. The LP has total opening
mass at most \(k\) and cost at most \(\OPT_k\).

\paragraph{Preparing a constant grid.}
Section~\ref{sec:preparation} replaces the LP openings by multiples of
\(g=1/\Delta\), where the even integer \(\Delta\) depends only on the
accuracy. The total mass increases by at most \(3g\), and each client's
expected assignment cost increases by an arbitrarily small factor.
Separated representative clients define disjoint bundles of nearby
opening mass. A bundle's deficit is the part missing from a unit; its
suppliers are the facilities or other bundles holding that part of the
representative's assignment. Rounding deficits together with their
suppliers controls both the opening budget and the cost of restoring a
full unit. This coupling must allow one supplier to serve many clients,
because assignment mass is not a shared capacity.

\paragraph{Rounding the grid.}
View every positive grid opening as copies of mass \(g\). In each round,
a directed graph specifies which copies a source erases when that source's
facility opens. Each target has exactly \(\Delta\) incoming sources,
chosen with distance-dependent probabilities. Opening and erasing are
balanced so that only a constant additional number of facilities is
needed. This source-opening and neighbor-erasure framework, including
balanced batches and forced large updates, builds on Byrka et
al.~\cite[Sections~3--6]{ByrkaEtAl2026}.

The connection-cost analysis needs a different comparison below factor
two. For each client, we track surviving copies from its initial assigned
unit and a persistent bound on its eventual connection distance. A
potential combines their remaining assignment cost with that bound. Its
coefficient depends on elapsed time as well as surviving mass. When one
of the client's own copies opens, we continue tracking the cheaper
survivors; this produces a subtraction that pays for other possible
erasures. The finite inequality in Section~\ref{sec:inequality} captures
that payment for arbitrary directed comparison coefficients. Its proof
transfers contributions from positive rows to lower-distance rows, whose
negative contributions cover them. Appendix~\ref{app:certificates}
proves every scalar estimate by explicit rational polynomial certificates.
The geometric substitution and the potential analysis are given in
Section~\ref{sec:rounding}.

A source erasing many copies is handled by forcing its facility open.
Because this choice depends on the graph, its potential change is bounded
for each realized graph first. Only then do we average over graphs. This
order lets the finite directed inequality control the averaged drift
without conditioning on a graph-dependent forcing event.

\paragraph{A deterministic solution and the exact budget.}
The preceding analysis uses randomness, but only specified low-order
moments of each local choice enter its bounds. Section~\ref{sec:derandomization}
constructs polynomial-size rational distributions preserving these
moments. All physical histories are expanded for a constant number of
rounds, and the cheapest output satisfying the additive budget is chosen.
The separate rational-arithmetic analysis proves polynomial time in the
binary input length. The result is a deterministic approximation opening
at most \(k+c\) facilities for an accuracy-dependent constant \(c\).
Li and Svensson's reduction then removes this surplus with an arbitrarily
small additional loss.

\paragraph{Organization.}
Section~\ref{sec:preliminaries} fixes the model and states the external
additive-to-ordinary reduction. Section~\ref{sec:preparation} prepares the
openings. Section~\ref{sec:inequality} proves the finite drift inequality,
with its scalar certificates in Appendix~\ref{app:certificates}.
Section~\ref{sec:rounding} gives and analyzes the iterative procedure.
Section~\ref{sec:derandomization} supplies the deterministic implementation
and completes the proof of Theorem~\ref{thm:main}.

